\answer{OA32}
维护一个双端队列，保存候选下标。队列内下标严格递增，对应数值严格递减。新值 $x_t$ 到来时依次执行：
\begin{enumerate}
\item 从队首删除所有小于 $t-N+1$ 的下标，即已离开窗口的元素。
\item 只要队列非空且队尾的值不大于 $x_t$，就删除队尾。
\item 将下标 $t$ 加入队尾，输出队首对应的值。
\end{enumerate}
删除队尾是安全的：若旧下标 $i<t$ 满足 $x_i\le x_t$，在任何同时包含它们的后续窗口内，$i$ 都不优于 $t$；而 $t$ 更晚过期。因此 $i$ 不会再成为必须保留的最大值候选。过期元素也不能影响当前窗口。更新后值按递减排列，队首就是当前最大值。相等值时保留较新的下标，因其有效期更长；若另需最早最大值下标，可把第二步的“不大于”改成“严格小于”。

每个下标只入队一次，至多出队一次。因此处理 $T$ 个值的总操作数为 $O(T)$，每次更新的摊还时间为 $O(1)$，队列空间为 $O(N)$。但一次很大的新值可以删除队列中几乎全部元素，所以单次最坏时间是 $O(N)$，不能把摊还界误写为逐次最坏界。$N=1$ 时输出就是新值本身。
